\section{Guide d'exécution}

Le projet peut être exécuté à partir de la ligne de commande avec :

\begin{lstlisting}[numbers=none]
python main.py
\end{lstlisting}

L'utilisateur sélectionne ensuite le problème, l'instance, la stratégie et le
mode de recherche souhaités.

\section{Dépôt de code}

Le code source complet, les instances de test et les scripts de benchmark sont
disponibles dans le dépôt suivant :

\begin{center}
\url{https://github.com/VOTRE-ORGANISATION/aussi-v2}
\end{center}

\section{Pseudocode du moteur générique}

\begin{lstlisting}[caption={Pseudocode conceptuel de la fonction SEARCH.}]
function SEARCH(problem, strategy, mode):

    root <- Node(problem.initial_state)
    OPEN <- frontier configured with strategy
    OPEN.push(root)

    best_g <- empty map
    if mode is GRAPH_SEARCH:
        best_g[key(root.state)] <- 0

    while OPEN is not empty:

        node <- OPEN.pop_best()

        if problem.is_goal(node.state):
            return build_success_result(node)

        for action in problem.actions(node.state):

            child_state <- problem.result(node.state, action)
            child_cost <- node.path_cost +
                          problem.step_cost(node.state, action, child_state)

            child <- Node(child_state, node, action,
                          child_cost, node.depth + 1)

            if mode is TREE_SEARCH:
                OPEN.push(child)

            else if child_cost improves best_g[child_state]:
                best_g[key(child_state)] <- child_cost
                OPEN.push(child)

    return failure_result()
\end{lstlisting}

\section{Répartition des contributions}

\begin{table}[H]
\centering
\caption{Répartition des contributions au sein de l'équipe.}
\begin{tabularx}{\textwidth}{p{3.5cm}Y}
\toprule
\textbf{Membre} & \textbf{Contributions principales} \\
\midrule
Nom 1 & Architecture générale, moteur de recherche, métriques. \\
Nom 2 & Labyrinthes, navigation, heuristiques. \\
Nom 3 & N-Puzzle, détection de solvabilité, limites de ressources. \\
Nom 4 & Jeu arithmétique, interface CLI, expériences et rapport. \\
\bottomrule
\end{tabularx}
\end{table}